% =====================================================================
% Główny plik raportu - Warsztaty AI
% "Miary perplexity a jakość agentów LLM"
% Klaudia Chwistek, 2026
%
% KOMPILACJA: w Overleafie ustaw Compiler na pdfLaTeX (Menu -> Settings
% -> Compiler). Ta preambuła wspiera pdfLaTeX z polskimi znakami.
% Jeśli wolisz XeLaTeX / LuaLaTeX, zakomentuj inputenc i odkomentuj
% fontspec (patrz niżej).
% =====================================================================

\documentclass[11pt,a4paper]{article}

% --- Polskie znaki + typografia ---
\usepackage[utf8]{inputenc}       % komentarz jeśli używasz XeLaTeX/LuaLaTeX
\usepackage[T1]{fontenc}
\usepackage{lmodern}
\usepackage[polish]{babel}

% Alternatywa dla XeLaTeX/LuaLaTeX (zakomentuj 3 linie wyżej):
% \usepackage{fontspec}
% \usepackage[polish]{babel}

% --- Matematyka ---
\usepackage{amsmath}
\usepackage{amssymb}

% --- Tabele i grafiki ---
\usepackage{graphicx}
\usepackage{booktabs}
\usepackage{float}

% --- Hyperlinks ---
\usepackage[hidelinks]{hyperref}

% --- Bibliografia ---
\usepackage[backend=biber,style=numeric,sorting=none]{biblatex}
\addbibresource{references.bib}

% --- Tytuł ---
\title{Miary perplexity a jakość agentów LLM\\
       \large Raport projektu --- Warsztaty AI}
\author{Klaudia Chwistek}
\date{Maj 2026}

\begin{document}

\maketitle

% =====================================================================
% Tutaj wkleisz potem inne sekcje - na razie tylko sekcja 2
% =====================================================================

% =====================================================================
% Sekcja 1 - Wprowadzenie
% =====================================================================

\section{Wprowadzenie}
\label{sec:wprowadzenie}

Perplexity (PPL) powstało jako miara jakości modelu językowego -
mówi, jak bardzo model jest zaskoczony obserwowanym tekstem. Z
czasem zaczęto je jednak traktować szerzej: jako tani sygnał o tym,
czy model dobrze radzi sobie z konkretnym zadaniem, liczony bez
potrzeby sprawdzania wyniku innym modelem czy człowiekiem. Pomysł
jest kuszący zwłaszcza w agentach LLM, gdzie taki sygnał mógłby
sterować decyzjami - kazać agentowi wstrzymać się z odpowiedzią,
poprosić o pomoc albo przeplanować działanie, gdy jego pewność spada.

Literatura nie jest tu jednak zgodna. Część prac pokazuje, że niższe
PPL wiąże się z lepszymi wynikami, inne - że PPL bywa zawodne i
podatne na artefakty długości czy interpunkcji. Większość tych badań
dotyczy przy tym pojedynczych, statycznych promptów. Tymczasem
agent nie produkuje jednej odpowiedzi, lecz całą trajektorię -
naprzemienne kroki rozumowania, wywołań narzędzi i obserwacji - i nie
jest oczywiste, czy wnioski z analizy pojedynczych promptów dają się
na taki przypadek przenieść.

Niniejsza praca bada to empirycznie. Sprawdzane jest, czy perplexity
i pokrewne sygnały rozkładowe (entropia tokenowa, top-1 margin)
przewidują poprawność odpowiedzi małego modelu LLM na zadaniach
matematycznych z GSM8K, w dwóch ustawieniach: czystego
chain-of-thought oraz agenta ReAct z kalkulatorem. Osobno badane
jest, w którym miejscu trajektorii agenta sygnał ten jest
najsilniejszy.

Z eksperymentów wyłaniają się trzy obserwacje. Metryki korzystające
z całego rozkładu prawdopodobieństw (entropia, margin) konsekwentnie
przewidują poprawność lepiej niż samo perplexity. Wewnątrz agenta
sygnał nie jest rozłożony równomiernie - najsilniej koreluje z
poprawnością pewność na segmentach formułujących wywołania narzędzia,
a praktycznie zanika na finalnej odpowiedzi. Wreszcie, choć pewność
porządkuje pytania sensownie, jako bezwzględna miara jest źle
wykalibrowana: model jest systematycznie zbyt pewny siebie. Ta ostatnia
obserwacja godzi sprzeczne stanowiska z literatury - PPL jest
użytecznym sygnałem porządkowym, ale zawodnym sygnałem ilościowym.

Dalsza część pracy ma następującą strukturę. Sekcja~\ref{sec:tlo}
omawia tło teoretyczne i powiązane prace. Sekcja~\ref{sec:metodologia}
opisuje zbiór danych, modele, sposób skoringu i badane metryki.
Sekcja~\ref{sec:wyniki} przedstawia wyniki korelacyjne w trzech
ustawieniach, a sekcja~\ref{sec:kalibracja} - analizę kalibracji.
Sekcja~\ref{sec:dyskusja} zestawia obserwacje z literaturą, a
sekcja~\ref{sec:limitations} omawia ograniczenia.

% =====================================================================
% Sekcja 2 - Tło teoretyczne (wersja kompaktowa, ~1 strona)
% =====================================================================

\section{Tło teoretyczne}
\label{sec:tlo}

\subsection{Perplexity jako miara modelu językowego}
\label{sec:perplexity-def}

Perplexity (PPL) to średnia geometryczna odwrotności prawdopodobieństw,
jakie model przypisuje kolejnym tokenom danej sekwencji. Dla
sekwencji $x_1, \dots, x_N$ pod modelem $p_\theta$:
\begin{equation}
\mathrm{PPL}(x_{1:N}) = \exp\!\left(-\frac{1}{N} \sum_{t=1}^{N} \log p_\theta(x_t \mid x_{<t})\right).
\label{eq:ppl}
\end{equation}
Niższa wartość oznacza, że model był mniej zaskoczony tym, co
zobaczył. Klasycznie PPL używa się jako miary jakości samego
modelu językowego - im niższe na zbiorze testowym, tym lepiej
model przewiduje tekst. W tej pracy interesuje jednak inne
zastosowanie: PPL liczone na własnej odpowiedzi modelu, jako
sygnał mówiący o tym, czy ta odpowiedź jest poprawna - taki,
żeby agent mógł np.\ odroczyć odpowiedź lub poprosić o pomoc,
gdy jego pewność spada.

\subsection{Argumenty za perplexity}
\label{sec:pro-ppl}

Gonen i in.~\cite{gonen2023demystifying} pokazali, że perplexity
samego promptu jest dobrym predyktorem skuteczności modelu - prompty
o niższym PPL dają wyższą trafność odpowiedzi. Interpretacja
autorów: prompty, które ,,brzmią naturalniej'' dla modelu (mniejsze
zaskoczenie na jego tokenach), są dla niego łatwiejsze i prowadzą
do lepszych wyników. Sam pomiar PPL działa tu jako tania heurystyka
wyboru lepszego promptu - bez konieczności faktycznego uruchamiania
zadania i sprawdzania wyniku.

\subsection{Argumenty przeciw perplexity}
\label{sec:anti-ppl}

Wang i in.~\cite{wang2022perplexity} pokazują drugą stronę medalu:
PPL bywa zawodne jako miara jakości tekstu. Autorzy zwracają uwagę
na trzy pułapki. Po pierwsze, krótsze teksty dostają wyższe PPL
niezależnie od jakości (\emph{length bias}). Po drugie, powtarzające
się fragmenty sztucznie je zaniżają (\emph{repetition sensitivity}).
Po trzecie, drobne różnice w interpunkcji potrafią wprowadzać dużą
zmienność (\emph{punctuation dependency}). Zestawienie tych dwóch
prac otwiera pytanie istotne dla tego raportu: czy ten obraz
zmienia się, gdy zamiast pojedynczego promptu analizowane są
wieloetapowe trajektorie agentów?

\subsection{Perplexity w długim kontekście}
\label{sec:longppl}

Fang i in.~\cite{fang2025longppl} pokazują problem klasycznego PPL
na zadaniach z długim kontekstem. Standardowe PPL uśrednia
logarytmy prawdopodobieństw po wszystkich tokenach z taką samą
wagą, przez co całość zdominowana jest przez tokeny lokalnie
przewidywalne - spójniki, końcówki fleksyjne, znaki interpunkcyjne -
a wkład tokenów, które naprawdę zależą od dalekiego kontekstu,
ginie w średniej. Zaproponowana przez autorów metryka LongPPL
waży tokeny zgodnie z tym, jak bardzo zależą od dalekiego kontekstu;
na benchmarkach długiego kontekstu osiąga ona korelację Pearsona
$r \approx -0{,}96$ z faktyczną wydajnością modeli - podczas gdy
zwykłe PPL daje na tych samych danych wartość bliską zeru.

Ten sam mechanizm może działać też w przypadku trajektorii agentów.
Trajektorie agenta ReAct - złożone z naprzemiennych bloków
\texttt{Thought}, \texttt{Action}, \texttt{Observation}
i \texttt{Answer} - są długie i pełne tokenów strukturalnych
(samych etykiet segmentów, znaków interpunkcyjnych), na których
model jest bardzo pewny niezależnie od jakości rozwiązania. Takie
tokeny mogą zdominować jedną wartość PPL liczoną na całej
trajektorii i wytłumić sygnał z miejsc, w których model
rzeczywiście ,,myśli'' - na przykład formułuje wywołanie
kalkulatora. Stąd w tej pracy oprócz PPL na całej trajektorii
liczone jest też PPL osobno na segmentach \texttt{Action}
i \texttt{Answer} (sprawdzane w sekcji~\ref{sec:wyniki-agent}).
Szerszy przegląd metryk ewaluacji agentów - w którym PPL pojawia
się obok trafności wywołań narzędzi, jakości planu, metryk
trajektorii i bezpieczeństwa - można znaleźć w pracy Mohammadi
i in.~\cite{mohammadi2025survey}.

\subsection{Niepewność jako sygnał sterujący w agentach}
\label{sec:agent-uncertainty}

Trzy nowsze artykuły idą krok dalej: zamiast tylko mierzyć pewność
modelu, wpinają ją wprost w pętlę decyzyjną agenta. ReDAct~\cite{redact2026}
wprowadza próg pewności, poniżej którego agent odracza odpowiedź
zamiast ryzykować błąd. \emph{Every Response Counts}~\cite{everyresponse2026}
zbiera niepewność wielu agentów pracujących nad tym samym zadaniem
i - korzystając z dekompozycji tensorowej - próbuje wyłowić, kiedy
modele systematycznie się ze sobą nie zgadzają. \emph{Towards
Agents That Know When They Don't Know}~\cite{towardsknow2025}
traktuje niską pewność jako sygnał, że trzeba się zatrzymać
i sprawdzić aktualny krok rozumowania, zanim agent pójdzie dalej.

Wszystkie trzy prace milcząco zakładają, że pewność liczona
z rozkładu tokenów rzeczywiście mówi coś o jakości tego, co model
generuje. W tej pracy sprawdzane jest, na ile to założenie się
utrzymuje na prostym agencie ReAct rozwiązującym zadania matematyczne:
jak silna jest korelacja PPL i entropii z poprawnością i czy ma
znaczenie, na której części trajektorii się ją mierzy.

% =====================================================================
% Sekcja 3 - Metodologia (wersja kompaktowa, ~1-1.5 strony)
% =====================================================================

\section{Metodologia}
\label{sec:metodologia}

\subsection{Pytanie badawcze}
\label{sec:rq}

W tej pracy sprawdzane jest, czy perplexity i pokrewne sygnały
liczone z rozkładu prawdopodobieństw tokenów - średnia entropia
tokenowa, top-1 margin, top-1 mass - przewidują poprawność
odpowiedzi LLM na konkretne pytanie. Pytanie to stawiane jest
w dwóch ustawieniach: w czystym chain-of-thought (CoT) oraz
w agencie ReAct z pojedynczym narzędziem.

Drugie, pomocnicze pytanie dotyczy tego, gdzie wewnątrz trajektorii
agenta najlepiej mierzyć tę pewność: czy mocniejszy sygnał płynie
z segmentów \emph{action} (wywołań narzędzia), \emph{thought}
(rozumowania) czy \emph{answer} (finalnej odpowiedzi).

\subsection{Zbiór danych}
\label{sec:gsm8k}

Eksperymenty przeprowadzono na zbiorze GSM8K~\cite{cobbe2021gsm8k} -
8500 zadań tekstowych z matematyki szkolnej (7473 trenujących
i 1319 testowych), z których każde wymaga od dwóch do ośmiu kroków
rozumowania. Każde zadanie kończy się jednoznaczną odpowiedzią
liczbową poprzedzoną sekwencją \texttt{\#\#\#\#}, dzięki czemu
ocena poprawności sprowadza się do porównania dwóch liczb - bez
potrzeby angażowania drugiego modelu jako sędziego. Przykład
z części testowej:

\begin{quote}
\textit{Natalia sprzedała spinacze 48 swoim przyjaciółkom w kwietniu,
a w maju sprzedała ich o połowę mniej. Ile spinaczy łącznie
sprzedała w kwietniu i maju?}\\
\textit{Rozwiązanie}: $48/2 = 24$ spinaczy w maju; $48 + 24 = 72$
spinaczy łącznie. \texttt{\#\#\#\# 72}
\end{quote}

GSM8K jest dla tego projektu wygodny z trzech powodów. Po pierwsze,
ocena jest zerojedynkowa i obiektywna - liczba albo się zgadza,
albo nie. Po drugie, zadania są krótkie i mieszczą się w oknie
kontekstowym małych modeli (0,5-1,5B parametrów). Po trzecie,
modele tej skali osiągają na GSM8K dokładność rzędu 20-60\%,
czyli ani nie zgadują na ślepo, ani nie rozwiązują wszystkiego -
co jest pożądane przy badaniu korelacji pewności z poprawnością.
GSM8K jest też standardowym benchmarkiem dla prac o rozumowaniu
krok po kroku od~\cite{wei2022cot}.

Z części testowej losowano 50 zadań (faza pilotażowa) lub 60 zadań
(eksperymenty główne) przy ustalonych ziarnach. Modele nie były
dotrenowywane - wszystkie wyniki pochodzą z czystej ewaluacji.

\subsection{Modele}
\label{sec:modele}

W eksperymentach użyto dwóch modeli z tej samej rodziny, różniących
się tylko skalą: Qwen2.5-0.5B-Instruct (w fazie pilotażowej) oraz
Qwen2.5-1.5B-Instruct (w eksperymentach głównych). Pozwala to
sprawdzić, czy obserwacje z fazy pilotażowej utrzymują się przy
przejściu na większy model. Spełnia to też wymóg co najmniej dwóch
modeli w ramach kursu Warsztaty AI.

\subsection{Dekodowanie i skoring}
\label{sec:dekodowanie}

We wszystkich eksperymentach stosowane jest dekodowanie zachłanne
(\emph{greedy decoding}) - na każdym kroku wybierany jest token
o najwyższym prawdopodobieństwie. Eliminuje to szum samplingu
i sprawia, że wartości PPL między pytaniami można porównywać.
Ma to swoją cenę: ponieważ model zawsze wybiera swój top-1 token,
jego prawdopodobieństwa są systematycznie wysokie, a wartości PPL
skupiają się blisko 1,0. Jest to jedno z ograniczeń, do których
raport wraca w sekcji~\ref{sec:limitations}.

Skoring działa w jednym przejściu modelu przez tekst złożony
z promptu i wygenerowanej kontynuacji. Tokeny promptu są maskowane,
więc do metryk wliczają się tylko te tokeny, które model sam
wygenerował - mierzona jest jego pewność co do własnego rozumowania,
a nie znajomość samego polecenia. Implementacja znajduje się w pliku
\texttt{experiments/perplexity.py} (funkcja \texttt{score\_continuation}).

\subsection{Dwa warianty: CoT i agent ReAct}
\label{sec:settings}

Pierwszy wariant to czysty chain-of-thought: model dostaje prompt
zachęcający do myślenia krok po kroku i generuje całą odpowiedź
w jednym przebiegu, bez narzędzi.

Drugi wariant to pętla ReAct~\cite{yao2023react}, w której model
na przemian produkuje segmenty oznaczone \texttt{Thought:},
\texttt{Action:}, \texttt{Observation:} i \texttt{Answer:}.
Jedynym dostępnym narzędziem jest kalkulator, a liczba wywołań
w jednym pytaniu ograniczona jest do pięciu.

Te same 60 pytań i ten sam model (Qwen2.5-1.5B) w obu wariantach
pozwalają czysto porównać, czy włożenie modelu w pętlę agentową
zmienia samą relację między PPL a poprawnością - nawet jeśli
dokładność spada lub rośnie.

\subsection{Metryki}
\label{sec:metryki}

Rozkład prawdopodobieństw tokenów sprowadzany jest do czterech
skalarów. Pierwszy z nich, \textit{perplexity}, dane jest
równaniem~\eqref{eq:ppl} i mówi, jak bardzo średnio model był
zaskoczony własnymi tokenami. \textit{Średnia entropia tokenowa}
$-\sum_v p(v)\log p(v)$, uśredniona po wszystkich pozycjach
sekwencji, korzysta z całego rozkładu, a nie tylko z tokena, który
został wybrany - dzięki temu reaguje także wtedy, gdy top-1 token
jest pewny, ale za nim stoi ,,gęsta'' alternatywa.
\textit{Top-1 margin} to różnica $p(\text{top-1}) - p(\text{top-2})$,
czyli to, jak bardzo wybrany token wygrał z najbliższym konkurentem.
\textit{Top-1 mass} to po prostu $p(\text{top-1})$.

Do oceny kalibracji wykorzystano \emph{Expected Calibration Error}
(ECE), liczone na binach o równej liczbie obserwacji
($n_{\text{bins}}=10$). Pewność modelu definiowana jest tu jako
$1/\mathrm{PPL}$. Niższe ECE oznacza lepszą zgodność tego, jak
pewny był model, z tym, jak często rzeczywiście miał rację.

% =====================================================================
% Sekcja 4 - Wyniki (wersja kompaktowa)
% Wartości liczbowe policzone bezpośrednio z plików w overnight_results/
% i results/ - są zsynchronizowane z notebookiem main_results_report.ipynb.
% =====================================================================

\section{Wyniki}
\label{sec:wyniki}

Wszystkie wyniki pochodzą z jednego przebiegu kodu z katalogu
\texttt{experiments/}, na ustalonych ziarnach losowości. Surowe
CSV z wynikami zapisane są w katalogu \texttt{overnight\_results/}
i można je odtworzyć z dołączonego notebooka
\texttt{main\_results\_report.ipynb}.

\subsection{Faza pilotażowa: Qwen-0,5B, czysty CoT, $n=150$}
\label{sec:wyniki-pilot}

Na 150 zadaniach (trzy ziarna po 50) model Qwen2.5-0,5B-Instruct
osiąga dokładność \textbf{44,7\%} - mieszczącą się w przedziale,
który dobrze nadaje się do badania korelacji pewności z poprawnością.

Już na tej skali widać pierwszy systematyczny wzorzec: średnia
entropia tokenowa jest lepszym predyktorem poprawności niż samo
perplexity ($r = -0{,}317$, $p < 0{,}001$ vs $r = -0{,}236$,
$p = 0{,}004$). Innymi słowy, sygnał z całego rozkładu
prawdopodobieństw niesie więcej informacji niż średnia liczona
tylko po tokenach, które model faktycznie wybrał. Ten sam wzorzec
powtarza się w każdym kolejnym ustawieniu.

\subsection{Agent ReAct z kalkulatorem: Qwen-1,5B, $n=60$}
\label{sec:wyniki-agent}

Na 60 zadaniach (dwa ziarna po 30) Qwen2.5-1,5B-Instruct
w pętli ReAct z pojedynczym narzędziem (kalkulator) osiąga
dokładność \textbf{28,3\%}. Spośród 60 trajektorii 12 (20\%)
zostaje obciętych - model albo przekracza limit pięciu wywołań
kalkulatora, albo łamie format dialogowy (np.\ zapomina o etykiecie
\texttt{Action:}).

Najciekawsza obserwacja w agencie dotyczy tego, \emph{gdzie} warto
mierzyć pewność. Metryki liczone tylko na segmentach \texttt{Action:} -
czyli tam, gdzie model formułuje wywołanie kalkulatora - dają
najsilniejszą korelację z poprawnością: \texttt{action\_ppl\_mean}
osiąga $r = -0{,}436$, $p = 0{,}002$, $n = 48$, z bootstrapowym
95\% CI $[-0{,}60,\ -0{,}24]$ (5000 resamplowań). Te same metryki
liczone tylko na finalnej odpowiedzi (\texttt{Answer:}) praktycznie
nie niosą żadnego sygnału ($r = -0{,}099$, $p = 0{,}50$, CI
$[-0{,}32,\ +0{,}21]$ - przedział ufności pokrywa zero).

Intuicja stojąca za tym wynikiem jest dość prosta. W momencie,
w którym model generuje finalną odpowiedź, wynik liczbowy znajduje
się już w jego kontekście - widoczny jako rezultat ostatniego
wywołania kalkulatora. Model w zasadzie kopiuje go z kontekstu,
więc jest bardzo pewny niezależnie od tego, czy całe rozumowanie
było prawidłowe. Pewność modelu w trakcie \emph{formułowania}
wywołań - gdy musi zdecydować, co policzyć - rozróżnia poprawne
i błędne odpowiedzi. Pewność na samej \emph{konkluzji} - już nie.

Sam rozkład pewności też różni się między typami segmentów.
Segmenty \texttt{Answer:} są skupione najmocniej (średnie PPL$=1{,}26$,
top-1 margin $=0{,}80$), segmenty \texttt{Action:} - pośrednio
(PPL$=1{,}46$, margin $=0{,}64$), segmenty \texttt{Thought:} - najmniej
(PPL$=1{,}59$, margin $=0{,}61$). Ten porządek utrzymuje się w obu
ziarnach.

\subsection{Kontrola czystego CoT: Qwen-1,5B, $n=60$}
\label{sec:wyniki-cot15}

Żeby zobaczyć, jak na te same pytania reaguje model bez pętli
agentowej, te same 60 zadań przepuszczono jeszcze raz przez
Qwen2.5-1,5B-Instruct, tym razem w czystym CoT. Dokładność rośnie
do \textbf{60,0\%} (wobec 28,3\% w wariancie agentowym), a najsilniejszy
sygnał korelacyjny w całym projekcie pojawia się właśnie tutaj:
średnia entropia tokenowa łańcucha rozumowania osiąga $r = -0{,}490$,
$p < 0{,}001$, $n = 60$. Samo perplexity też koreluje z poprawnością,
ale słabiej: $r = -0{,}283$, $p = 0{,}028$. Włożenie modelu w pętlę
ReAct nie likwiduje więc sygnału PPL/entropii - przenosi go z całej
kontynuacji na segmenty wywołań narzędzia.

\subsection{Porównanie zbiorcze}
\label{sec:headline}

Tabela~\ref{tab:headline} zestawia korelacje Pearsona między metrykami
a poprawnością w trzech głównych ustawieniach.

\begin{table}[h]
\centering
\small
\begin{tabular}{llrrr}
\toprule
\textbf{Ustawienie} & \textbf{Metryka} & $n$ & $r$ & $p$ \\
\midrule
0,5B CoT (pilot)        & \texttt{cot\_perplexity}           & 150 & $-0{,}236$ & $0{,}004$ \\
0,5B CoT (pilot)        & \texttt{cot\_mean\_token\_entropy} & 150 & $-0{,}317$ & $<0{,}001$ \\
1,5B CoT                & \texttt{cot\_perplexity}           & 60  & $-0{,}283$ & $0{,}028$ \\
1,5B CoT                & \texttt{cot\_mean\_token\_entropy} & 60  & $-0{,}490$ & $<0{,}001$ \\
1,5B Agent (cała traj.) & \texttt{agent\_ppl\_mean}          & 60  & $-0{,}354$ & $0{,}006$ \\
1,5B Agent (cała traj.) & \texttt{agent\_entropy\_mean}      & 60  & $-0{,}388$ & $0{,}002$ \\
1,5B Agent (action)     & \texttt{action\_ppl\_mean}         & 48  & $-0{,}436$ & $0{,}002$ \\
1,5B Agent (action)     & \texttt{action\_margin\_mean}      & 48  & $+0{,}454$ & $0{,}001$ \\
\bottomrule
\end{tabular}
\caption{Korelacje Pearsona między metrykami a poprawnością
w trzech ustawieniach ($n$ to liczba obserwacji po usunięciu
braków danych). Wartości policzone bezpośrednio z plików CSV
w katalogu \texttt{overnight\_results/}.}
\label{tab:headline}
\end{table}

Z tabeli wyłaniają się dwa powtarzające się wzorce, niezależne
od skali modelu i ustawienia. Po pierwsze, średnia entropia tokenowa
i top-1 margin konsekwentnie wypadają lepiej niż samo perplexity
jako predyktory poprawności - pełny rozkład prawdopodobieństw
niesie więcej informacji niż prawdopodobieństwa tylko wybranych
tokenów. Po drugie, wewnątrz agenta liczy się to, na której części
trajektorii mierzy się pewność: korelacje na segmentach
\texttt{action} są silniejsze niż na całych trajektoriach, a te
z kolei silniejsze niż na samym segmencie \texttt{answer}.

% =====================================================================
% Sekcja 5 - Kalibracja
% Wartości z overnight_results/RESULTS_CALIBRATION_15B.md (calibration_15b.png).
% =====================================================================

\section{Kalibracja}
\label{sec:kalibracja}

Wyniki z sekcji~\ref{sec:wyniki} pokazują, że pewność modelu
koreluje z poprawnością. Korelacja jest jednak własnością czysto
porządkową - mówi tylko tyle, że pytania da się ustawić od
najbardziej do najmniej ryzykownego. Nie odpowiada natomiast na
pytanie, czy pewność rzędu 0,85 faktycznie oznacza 85\% szans na
poprawną odpowiedź. To już kwestia kalibracji i właśnie jej
dotyczy ta sekcja.

Jako pewność modelu przyjęto $1/\mathrm{PPL}$ - liczbę z przedziału
$(0, 1]$. Jej zgodność z poprawnością mierzy Expected Calibration
Error (ECE) na binach o równej liczbie obserwacji
($n_{\text{bins}}=10$). Biny o równej liczności są tu potrzebne,
bo przy dekodowaniu zachłannym wartości $1/\mathrm{PPL}$ skupiają
się blisko 1,0 (sekcja~\ref{sec:dekodowanie}) i biny o stałej
szerokości byłyby w większości puste.

\begin{table}[h]
\centering
\small
\begin{tabular}{lrrr}
\toprule
\textbf{Ustawienie} & $n$ & \textbf{dokładność} & \textbf{ECE} \\
\midrule
1,5B CoT                       & 60 & 60,0\% & $0{,}265$ \\
1,5B Agent (cała traj.)        & 60 & 28,3\% & $0{,}425$ \\
1,5B Agent (segmenty action)   & 48 & 25,0\% & $0{,}445$ \\
\bottomrule
\end{tabular}
\caption{Expected Calibration Error dla pewności zdefiniowanej
jako $1/\mathrm{PPL}$ (binowanie o równej liczności,
$n_{\text{bins}}=10$). Niższe ECE oznacza lepszą zgodność
deklarowanej pewności z faktyczną poprawnością.}
\label{tab:ece}
\end{table}

ECE jest wysokie w każdym ustawieniu (tabela~\ref{tab:ece}).
Diagramy niezawodności (plik \texttt{calibration\_15b.png})
tłumaczą skąd to się bierze: w prawie każdym binie średnia pewność
leży powyżej średniej poprawności, czyli model jest przesadnie
pewny. W wariancie agentowym widać to najmocniej - w najniższych
binach pewności model deklaruje 0,49-0,66, podczas gdy poprawność
w tych binach wynosi 0\%. Nawet tam, gdzie agent jest najmniej
pewny, i tak jest pewniejszy, niż uzasadniałyby to jego wyniki.

Zestawienie tej obserwacji z sekcją~\ref{sec:wyniki} prowadzi do
głównego wniosku. Perplexity liczone na rozumowaniu modelu istotnie
koreluje z poprawnością, a mimo to daje wysokie ECE. Dobry porządek
nie pociąga więc za sobą dobrej kalibracji - jedno z drugiego nie
wynika.

Ma to konkretne znaczenie dla zastosowań. Tam, gdzie liczy się
sama kolejność, perplexity się sprawdza: można nim wybrać
najpewniejszą z kilku odpowiedzi albo wskazać najbardziej ryzykowne
pytania do ręcznego sprawdzenia. Gorzej z regułami opartymi na
stałym progu, w stylu ,,odrocz odpowiedź, gdy pewność spadnie
poniżej 0,7'' - wartość 0,7 nie oznacza tu 70\% szans na poprawność,
lecz znacznie mniej, więc taki próg trzeba najpierw wykalibrować
na danych z konkretnego zadania. Jest to zastrzeżenie wobec prac
omawianych w sekcji~\ref{sec:agent-uncertainty}: kierunek sygnału
jest dobry, ale jego skala nie nadaje się do bezpośredniego progowania.

% =====================================================================
% Sekcja 6 - Dyskusja
% =====================================================================

\section{Dyskusja}
\label{sec:dyskusja}

Zebrane wyniki układają się w kilka spójnych wątków, które warto
omówić razem - bo dopiero zestawione ze sobą i z literaturą z
sekcji~\ref{sec:tlo} mówią coś więcej niż pojedyncze współczynniki
korelacji.

Pierwszy wątek dotyczy wyboru metryki. We wszystkich ustawieniach -
od pilotażowego modelu 0,5B po agenta 1,5B - średnia entropia
tokenowa i top-1 margin okazują się lepszymi predyktorami
poprawności niż samo perplexity. Różnica jest niewielka, ale
konsekwentna, i ma dość naturalne wytłumaczenie: entropia i margin
korzystają z całego rozkładu prawdopodobieństw, a perplexity tylko
z prawdopodobieństw tokenów, które model faktycznie wybrał. Model
może wybrać token z wysokim prawdopodobieństwem, mając tuż za nim
silną alternatywę - perplexity tego nie zauważy, a margin i entropia
tak. Sugeruje to, że w dalszych pracach nad pewnością agentów warto
sięgać po metryki rozkładowe, a nie ograniczać się do klasycznego PPL.

Drugi wątek to lokalizacja pomiaru wewnątrz agenta. Sygnał nie jest
rozłożony równomiernie po trajektorii: najsilniej koreluje z
poprawnością pewność na segmentach \texttt{Action}, słabiej na całej
trajektorii, a na samym segmencie \texttt{Answer} znika niemal
zupełnie. Wynika to z mechaniki ReAct - w momencie formułowania
finalnej odpowiedzi wynik liczbowy jest już w kontekście, więc model
go po prostu przepisuje i jest pewny niezależnie od tego, czy droga
do tego wyniku była poprawna. Pewność niesie informację tam, gdzie
model rzeczywiście podejmuje decyzję (co policzyć), a nie tam, gdzie
tylko odtwarza gotowy rezultat. To ustalenie jest zbieżne z
obserwacją Fanga i in.~\cite{fang2025longppl}, że uśrednianie po
całej długiej sekwencji rozcieńcza sygnał - tyle że tutaj rolę
,,nieistotnych'' tokenów grają nie spójniki, lecz cała końcowa
faza kopiowania odpowiedzi.

Trzeci, najważniejszy wątek dotyczy różnicy między korelacją a
kalibracją (sekcja~\ref{sec:kalibracja}). Perplexity porządkuje
pytania sensownie, ale jako bezwzględna miara pewności jest źle
wykalibrowane - model jest systematycznie zbyt pewny siebie, i to
tym bardziej, im trudniejsze ustawienie. Rozstrzyga to pozornie
sprzeczne stanowiska z literatury. Gonen i
in.~\cite{gonen2023demystifying} mają rację, że niższe PPL wiąże się
z lepszymi wynikami - kierunek jest poprawny. Zastrzeżenia Wanga i
in.~\cite{wang2022perplexity} co do wiarygodności PPL też są jednak
zasadne, jeśli traktować je jako uwagę o skali, a nie o porządku.
Te dwa stanowiska nie są sprzeczne - opisują dwie różne własności
tej samej liczby.

Płynie stąd praktyczny wniosek dla prac, które wbudowują pewność w
pętlę decyzyjną agenta (sekcja~\ref{sec:agent-uncertainty}). ReDAct
\cite{redact2026} i pokrewne podejścia opierają się na progowaniu
surowej pewności, a wyniki tej pracy pokazują, że taki próg nie
przenosi się między zadaniami bez rekalibracji: ta sama wartość
liczbowa oznacza inną faktyczną szansę poprawności w czystym CoT i
w agencie. Sam pomysł sterowania niepewnością pozostaje sensowny,
ale wymaga albo kalibracji na danych docelowych, albo oparcia decyzji
na względnym uporządkowaniu, a nie na sztywnym progu.

Na koniec warto odnotować, że dwa główne wzorce - przewaga metryk
rozkładowych oraz przeniesienie sygnału na segmenty \texttt{Action} -
utrzymują się przy przejściu z modelu 0,5B na 1,5B oraz między
ziarnami losowości. Przy skali tego projektu nie jest to dowód
ogólności, ale sugeruje, że obserwacje nie są wyłącznie artefaktem
jednego modelu czy jednego losowania danych.

% =====================================================================
% Sekcja 7 - Ograniczenia
% =====================================================================

\section{Ograniczenia}
\label{sec:limitations}

Wyniki tej pracy trzeba czytać z kilkoma zastrzeżeniami.

Najpoważniejsze dotyczy dekodowania. Wszystkie eksperymenty
korzystają z dekodowania zachłannego, przez co model na każdym kroku
wybiera swój najbardziej prawdopodobny token, a wartości
$1/\mathrm{PPL}$ skupiają się blisko 1,0. Sygnał pewności jest więc
ściśnięty i część informacji, którą niosłoby próbkowanie z pełnego
rozkładu, zostaje utracona. Powtórzenie eksperymentów z włączonym
samplowaniem mogłoby dać bogatszy obraz - i prawdopodobnie inne,
być może niższe ECE.

Ograniczona jest też skala. Badano dwa małe modele z jednej rodziny
(Qwen2.5 0,5B i 1,5B) na jednym zbiorze - GSM8K. To zadania
arytmetyczne z dokładnym dopasowaniem odpowiedzi liczbowej, więc nie
wiadomo, na ile obserwacje przenoszą się na większe modele ani na
zadania o innym charakterze (otwarte odpowiedzi, kod, dłuższe
trajektorie z wieloma narzędziami). Agent korzystał z jednego
narzędzia - kalkulatora - co jest najprostszym możliwym ustawieniem
ReAct.

Niewielkie są wreszcie próby. Korelacje liczone są na 60 (a w pilocie
150) obserwacjach, a w analizie segmentów \texttt{Action} efektywne
$n$ spada do 48. Przy takich liczebnościach przedziały ufności są
szerokie - bootstrapowy 95\% CI dla najmocniejszej korelacji sięga
od $-0{,}60$ do $-0{,}24$ - więc raportowane współczynniki należy
traktować jako rzędy wielkości i kierunki zależności, a nie jako
precyzyjne oszacowania. Z tego samego powodu w pracy nie testowano
rekalibracji pewności; wskazano ją jedynie jako naturalny dalszy krok.


% --- Bibliografia ---
\printbibliography

\end{document}
